\documentclass[11pt,final]{amsart}

\usepackage{amsmath,amssymb,amsthm,mathtools,mathrsfs}
\usepackage[margin=1in]{geometry}

\renewcommand{\vec}[1]{\boldsymbol{#1}}
\newcommand{\defeq}{\vcentcolon=}
\newcommand{\mat}[1]{\mathsf{\mathbf{#1}}}
\newcommand{\GM}[2]{\mathcal{N}\!\left( {#1}, {#2}\right)}
\newcommand{\obs}{{y}} 
\newcommand{\prcov}{\Gamma_{\mathrm{pr}}}
\newcommand{\postcov}{\Gamma_{\mathrm{post}}}
\newcommand{\noise}{\Gamma_{\mathrm{noise}}}
\newcommand{\R}{\mathbb{R}}
\newcommand{\B}{\mathrm{B}}
\newcommand{\F}{\mathrm{F}}
\newcommand{\I}{\mathrm{I}}
\newcommand{\St}{\widetilde{\mathrm{S}}}
\newcommand{\dpar}{{m}}
\newcommand{\dparpr}{\dpar_{\mathrm{pr}}}
\newcommand{\dparpost}{\dpar_{\mathrm{post}}}
\newcommand{\priorm}{\mu_{\mathrm{pr}}}
\newcommand{\postm}{\mu_{\mathrm{post}}}
\newcommand{\priorml}{\mu_{\mathrm{pr}}^\ell}
\newcommand{\postml}{\mu_{\mathrm{post}}^\ell}
\newcommand{\tran}{{\mkern-1.5mu\mathsf{T}}}
\newcommand{\ip}[2]{\langle {#1}, {#2}\rangle}
\newcommand{\lprmean}{\ell_\text{pr}^\text{mean}}
\newcommand{\lprsig}{\ell_\text{pr}^\text{var}}
\newcommand{\lpostmean}{\ell_\text{post}^\text{mean}}
\newcommand{\lpostsig}{\ell_\text{post}^\text{var}}
%\newcommand{\Gammay}{{\Gamma}_{y}}
\newcommand{\Gammay}{\Sigma}
\newcommand{\hilb}{\mathcal{H}}
\newcommand{\M}{\mathcal{M}}
\newcommand{\RR}{\mathcal{R}}
\newcommand{\LL}{\mathfrak{L}}
\DeclareMathOperator{\Cov}{Cov}
\DeclareMathOperator{\Var}{Var}
\newcommand{\CM}{\mathcal{E}}

\theoremstyle{plain}
\newtheorem{theorem}{Theorem}[section]
\newtheorem{proposition}[theorem]{Proposition}
\newtheorem{lemma}[theorem]{Lemma}
\newtheorem{corollary}[theorem]{Corollary}
\newtheorem{remark}[theorem]{Remark}
\newtheorem{definition}[theorem]{Definition}

\title{Some questions of learnability in linear inverse problems}
\author{Alen Alexanderian and Ruanui Nicholson}
\date{\today}

\begin{document}

\maketitle

\begin{abstract}
Insert abstract here.
\end{abstract}

\section{Introduction}
Let $\M$ be an infinite dimensional real separable Hilbert space. 
We consider the inverse problem of estimating $m \in \M$ from 
\begin{equation}\label{equ:model}
y = \F m + \eta
\end{equation}
We organize the present study along the following directions:
\begin{itemize}
\item \textbf{Informativeness of the model}: 
Given a fixed prior, to what extent can we improve the learnability of $m$ by 
choosing an optimal model $\F$? 

\item \textbf{Learnability of linear prediction functionals}:
Given 
a fixed prior and model $\F$, which linear prediction functionals of $m$ 
are most learnable? 
\end{itemize}

Alen: let's refine these questions. I have core-dumped both of the write-ups in the 
following sections. 

\section{Preliminaries}

Let $\M$ be an infinite-dimensional real separable Hilbert space, and let
$\F\in \mathcal{L}(\M,\R^q)$.
Consider the linear inverse problem of estimating $\dpar \in \M$ from the observation 
model
\[
    \obs = \F \dpar + {\eta},
\]
where ${\eta} \sim \GM{0}{\noise}$ and $\noise\in\R^{q\times q}$ is symmetric
positive definite. We assume a Gaussian prior 
$\GM{\dparpr}{\prcov}$.
Herein, $\dparpr\in\M$ and 
$\prcov$ is a strictly positive selfadjoint trace class operator on $\M$.
The posterior distribution is Gaussian 
$\postm = \GM{\dparpost}{\postcov}$. 
In particular, 
letting 
\begin{equation}\label{equ:Gammay}
\Sigma \defeq \F\prcov\F^*+\noise,
\end{equation} 
the posterior mean and covariance can be represented as 
\[
\begin{aligned}
    \dparpost &= \dparpr + \prcov \F^* \Gammay^{-1} (\obs-\F\dparpr),
    \\
    \postcov &= \prcov - \prcov \F^* \Gammay^{-1} \F\prcov.
\end{aligned}
\]
We consider posterior uncertainty in bounded linear functionals 
of the inversion parameter. 
In the analysis that follows, it will be convenient to work 
with the associated Riesz maps:
\begin{definition}
Let $\hilb$ be a real Hilbert space with inner product
$\ip{\cdot}{\cdot}_{\hilb}$.
The \emph{Riesz map} of $\hilb$ is the mapping
$\RR_{\hilb}:\hilb\to\hilb'$
defined by
\[
\bigl(\RR_{\hilb}v\bigr)(x)
=
\ip{v}{x}_{\hilb},
\qquad \text{for all } v,x\in\hilb.
\]
The Riesz map is a linear isometric isomorphism.
Further, by the Riesz Representation Theorem,
every $\ell\in\hilb'$ admits a unique representation
\[
\ell(x)=\ip{v_\ell}{x}_{\hilb},
\qquad \text{for all } x\in\hilb,
\]
with $v_\ell=\RR_{\hilb}^{-1}\ell$.
\end{definition}

Consider a bounded linear functional $\ell$ on $\M$ and let 
$v_\ell \defeq \RR_\hilb^{-1}\ell$ be its Riesz representer,
\begin{equation}\label{equ:functional}
    \ell(\dpar)
    =
    \ip{v_\ell}{\dpar}.
\end{equation}
Denote the posterior and prior laws of $\ell$ by
$\priorml$ and $\postml$, respectively.
Since $\ell$ is linear, these measures are both Gaussian: 
$\priorml = \GM{\lprmean}{\lprsig}$ and $\postml = \GM{\lpostmean}{\lpostsig}$ 
with 
\begin{subequations}\label{equ:moments}
\begin{align}
\lprmean &= \ell(\dparpr),\\
\lprsig &= \ip{v_\ell}{\prcov v_\ell},\\
    \lpostmean 
        &= \ell(\dparpr) + \ip{\F\prcov v_\ell}{\Gammay^{-1} (\obs-\F\dparpr)},\\
    \lpostsig 
        &= \ip{\prcov  v_\ell}{v_\ell} - \ip{\F\prcov v_\ell}{\Gammay^{-1}(\F\prcov v_\ell)}.
\end{align}
\end{subequations}

\section{Informativeness of the forward model}
\subsection{Infinite dimensions}

We consider optimal experimental design (OED) in the Bayesian setting for infinite dimensional linear inverse problems. In particular, let $m$ have Gaussian prior law $\mathcal{N}(m_0,\mathcal{C})$ on the Hilbert space $\mathcal{H}$. The prior covariance operator $\mathcal{C}$ is assumed to be {\em strictly} positive self-adjoint and trace class. In what follows it will useful to consider the spectral decomposition of the prior covariance operator,
\begin{align}
    \mathcal{C}=\sum_{i=1}^\infty\lambda_i\phi_i\otimes\phi_i,\label{eq:KLexp}
\end{align}
where $\lambda_1\geq \lambda_2\geq\cdots$ with $\lambda_i>0$ for $i\in\{1,2,\dots\}.$ We will assume to begin with that the multiplicity of $\lambda_1$ is 1.

We initially consider the set up 
\[y=\ell(m)+e=\langle g,m\rangle_\mathcal{H} + e=Am+e,\]
where $A\in\mathcal{L}(\mathcal{H},\mathbb{R})$ and $e\sim\mathcal{N}(0,\sigma^2_e).$ Then it is well known (see e.g., \cite[Equation (3.4b)]{stuart2010inverse}) that the posterior covariance operator is of the form
\[\mathcal{C}_{\rm post}=\mathcal{C}-\mathcal{C}A^\ast(\sigma_e^2+A\mathcal{C}A^\ast)^{-1}A\mathcal{C}\].

We are interested in the {\em direction} to measure not the {\em strength}, so we restrict ourself to the setting $AA^\ast=\|A\|_{\mathcal{H}^\ast}=1$. This (Reisz) implies that we have $\|g\|_\mathcal{H}=\|A\|_{\mathcal{H}^\ast}=1$

Now, by assumption $\mathcal{C}$ is trace class, and thus we can consider the difference
\begin{align}
    \text{Tr}(\mathcal{C})-\text{Tr}(\mathcal{C}_{\rm post})&=\text{Tr}(\mathcal{C}A^\ast(\sigma_e^2+A\mathcal{C}A^\ast)^{-1}A\mathcal{C})\nonumber\\
    &=\frac{1}{(\sigma_e^2+A\mathcal{C}A^\ast)}\text{Tr}(\mathcal{C}A^\ast A\mathcal{C})\nonumber\\
    &=\frac{1}{(\sigma_e^2+A\mathcal{C}A^\ast)}\text{Tr}( A\mathcal{C}\mathcal{C}A^\ast)\nonumber\\
    &=\frac{1}{(\sigma_e^2+A\mathcal{C}A^\ast)}  A\mathcal{C}\mathcal{C}A^\ast.\nonumber
\end{align}
Using the spectral decomposition (\ref{eq:KLexp}), we have
\begin{align}
    \text{Tr}(\mathcal{C})-\text{Tr}(\mathcal{C}_{\rm post})&=\frac{1}{(\sigma_e^2+A\sum_{i=1}^\infty\lambda_i\phi_i\otimes\phi_iA^\ast)} A\sum_{i=1}^\infty\lambda_i\phi_i\otimes\phi_i \sum_{j=1}^\infty\lambda_j\phi_j\otimes\phi_j A^\ast\nonumber\\
    &=\frac{1}{(\sigma_e^2+A\sum_{i=1}^\infty\lambda_i\phi_i\otimes\phi_iA^\ast)} A\sum_{i=1}^\infty\lambda_i^2\phi_i\otimes\phi_i  A^\ast\nonumber\\
    &=\frac{1}{(\sigma_e^2+\sum_{i=1}^\infty\lambda_i(A\phi_i)^2)}\sum_{i=1}^\infty\lambda_i^2(A\phi_i)^2.\nonumber\\
    &\leq \frac{\lambda_1}{(\sigma_e^2+\sum_{i=1}^\infty\lambda_i(A\phi_i)^2)}\sum_{i=1}^\infty\lambda_i(A\phi_i)^2,
    \nonumber\\
    &\leq \frac{\lambda_1^2}{(\sigma_e^2+\lambda_1\sum_{i=1}^\infty(A\phi_i)^2)}\sum_{i=1}^\infty(A\phi_i)^2,\label{eq:diff}
\end{align}
since for $a>0$ the function $\frac{x}{a+x}$ is increasing in $x$ and by assumption $\lambda_1\geq \lambda_2\geq\cdots$ with $\lambda_i>0$ for $i\in\{1,2,\dots\}.$ Lastly, using the fact that $\{\phi_i\}_{i=1}^{\infty}$ forms an orthonormal basis of $\mathcal{H}$, we have 
\[A\phi_i=\langle g,\phi_i\rangle_\mathcal{H}=\left\langle \sum_{j=1}^{\infty}\langle\phi_j,g\rangle\phi_j,\phi_i\right\rangle_\mathcal{H},\]
meaning that we have equality in (\ref{eq:diff}) when $A\phi_i=\langle \phi_1,\phi_i\rangle_\mathcal{H}$ for $i=1,2,\dots$. That is to say, the greatest decrease in the trace is when $g=\phi_1$ and is given by
\begin{align}
    \text{Tr}(\mathcal{C})-\text{Tr}(\mathcal{C}_{\rm post})&= \frac{1}{(\sigma_e^2+\lambda_1\sum_{i=1}^\infty(A\phi_i)^2)}\lambda_1^2\sum_{i=1}^\infty(A\phi_i)^2\nonumber\\
    &= \frac{\lambda_1^2}{\sigma_e^2+\lambda_1}\nonumber
\end{align}



\subsection{Finite dimensions}
Suppose we are interested in a parameter $\mathbb{R}^n\ni m\sim\pi(m)=\mathcal{N}(m_0,\Gamma_m)$. We {\em wish} to collect data $d\in\mathbb{R}$ related through a linear mapping of the form
\[d=w^\top m+e,\quad w\in\mathbb{R}^n,\quad \mathbb{R}\ni e\sim\mathcal{N}(0,\sigma_e^2).\]
Then what is the {\em best} $w\in\mathbb{R}^n$ to choose. Of course `best' is up to interpretation, but if go with reducing the uncertainty the most, then A- and D-optimality are natural metrics. In any case, the posterior covariance matrix is of the form
\begin{align}
   \Gamma_{\rm post}=(w\sigma_e^{-2}w^T+\Gamma_m^{-1})^{-1}=\Gamma_m-\Gamma_m w(w^\top\Gamma_m w+\sigma_e^2)^{-1}w^\top\Gamma_m.\label{eq:post} 
\end{align}


\paragraph{A-optimality}
For A-optimality, the reduction in the trace from prior to posterior is given by
\begin{align}
    \text{tr}(\Gamma_m)-\text{tr}(\Gamma_{\rm post})&=\text{tr}(\Gamma_m w(w^\top\Gamma_m w+\sigma_e^2)^{-1}w^\top\Gamma_m)\nonumber\\
    &=\text{tr}((w^\top\Gamma_m w+\sigma_e^2)^{-1}w^\top\Gamma_m\Gamma_m w)\nonumber\\
    &=\frac{w^\top\Gamma_m\Gamma_m w}{w^\top\Gamma_m w+\sigma_e^2}.\label{eq:trace}
\end{align}



Now, suppose $\Gamma_m$ is a {\em proper} covariance matrix, so that we can decompose it as
\begin{align}
    \Gamma_m=V\Lambda V^\top, \label{eq:PCA}
\end{align}
where $\mathbb{R}^{n \times n}\ni V=[v_1,v_2,\dots,v_n]$, $\mathbb{R}^{\times n}\ni\Lambda=\text{diag}(\lambda_1,\lambda_2,\dots,\Lambda_n)$ with  $\lambda_1\geq\lambda_2\geq\cdots\geq\lambda_n>0$ and
\[\Gamma_m v_i=\lambda_i v_i,\quad i=1,2,\dots n.\]
Then the reduction in the trace (\ref{eq:trace}) is given by
\begin{align}
    \text{tr}(\Gamma_m)-\text{tr}(\Gamma_{\rm post})&=\frac{w^\top\Gamma_m\Gamma_m w}{w^\top\Gamma_m w+\sigma_e^2}=\frac{w^\top V\Lambda V^\top V\Lambda V^\top w}{w^\top V\Lambda V^\top w+\sigma_e^2}\nonumber\\
    &=\frac{w^\top V\Lambda^2 V^\top w}{w^\top V\Lambda V^\top w+\sigma_e^2}.\label{eq:trace2}
\end{align}

Any $w$ can be written in terms of $V$, i.e., 
\[w=V\beta,\quad \beta\in\mathbb{R}^n.\]
We are only interested in the {\em direction} of the measurement, so we can restrict our attention to $\|w\|_2=\|\beta\|_2=1.$ Now, plugging this into (\ref{eq:trace2}), we have
\begin{align}
    \text{tr}(\Gamma_m)-\text{tr}(\Gamma_{\rm post})&=\frac{w^\top V\Lambda^2 V^\top w}{w^\top V\Lambda V^\top w+\sigma_e^2}=\frac{\beta^\top V^\top V\Lambda^2 V^\top V \beta}{\beta^\top V^\top V\Lambda V^\top  V \beta+\sigma_e^2}\nonumber\\
    &=\frac{\beta^\top \Lambda^2 \beta}{\beta^\top \Lambda\beta+\sigma_e^2}=\frac{\sum_{i=1}^n\lambda_i^2\beta_i^2}{\sum_{i=1}^n\lambda_i\beta_i^2+\sigma_e^2}\nonumber.
\end{align}
Since $\beta_i^2\geq0$ and $\lambda_1\geq\lambda_2\geq\cdots\geq\lambda_n>0$, we have $\lambda_1\sum_{i=1}^n\lambda_i\beta_i^2\geq\sum_{i=1}^n\lambda_i^2\beta_i^2$, so that
\begin{align}
    \text{tr}(\Gamma_m)-\text{tr}(\Gamma_{\rm post})&=\frac{\sum_{i=1}^n\lambda_i^2\beta_i^2}{\sum_{i=1}^n\lambda_i\beta_i^2+\sigma_e^2}\leq\frac{\lambda_1\sum_{i=1}^n\lambda_i\beta_i^2}{\sum_{i=1}^n\lambda_i\beta_i^2+\sigma_e^2}\leq\frac{\lambda_1^2}{\lambda_1+\sigma_e^2},\nonumber
\end{align}
where the last inequality follows (again) from the fact that $\sum_{i=1}^n\lambda_i\beta_i^2\leq\lambda_1\sum_{i=1}^n\beta_i^2=\lambda_1$, and that $\frac{\lambda_1\sum_{i=1}^n\lambda_i\beta_i^2}{\sum_{i=1}^n\lambda_i\beta_i^2+\sigma_e^2}$ is increasing in $\lambda$. The last step is to see that equality holds if $\beta=e_1=[1,0,0,\dots,0]^T\in\mathbb{R}^n$. That is to say, setting $w=v_1$ reduces the the uncertainty the most (using A-opt).

I'm thinking that if $\lambda_1=\lambda_2=\cdots=\lambda_r>\lambda_{r+1}\geq\lambda_{r+2}\geq\cdots\geq\lambda_n>0$, i.e., the largest eigenvalue has multiplicity greater than 1, then the result is something like take $w$ to be the eigenspace. But actually haven't sorted that out yet.

\paragraph{D-optimality}
For D-optimality (finite dimensional case), we will use  the{\em Matrix determinant lemma}: If $A\in\mathbb{R}^{n\times n}$ be invertible and $u,v\in\mathbb{R}^n$ then
    \[\det{(A+uv^\top)}=\det{(1+v^\top A^{-1} u)}\det{(A)},\] 
 and the fact that $\det(A^{-1})=\det(A)^{-1}$.
 First,
\begin{align}
   \det{(\Gamma_{\rm post})}&=\det{((w\sigma_e^{-2}w^T+\Gamma_m^{-1})^{-1})}\nonumber\\
&=\det{(w\sigma_e^{-2}w^T+\Gamma_m^{-1})}^{-1}\nonumber\\
   &=((1+\sigma_e^{-2}w^\top\Gamma_m w)\det{(\Gamma_m^{-1})})^{-1}\nonumber\\
   &=\frac{\det{(\Gamma_m)}}{(1+\sigma_e^{-2}w^\top\Gamma_m w)}\label{eq:detpost}
\end{align}
The key quantity for D-optimality is 
\[\log{\left(\frac{\det{(\Gamma_m)}}{\det{(\Gamma_{\rm post})}}\right)}=\log\det{(\Gamma_m)}-\log\det{(\Gamma_{\rm post})}.\]
So, looking at the log-det of the posterior covariance matrix (\ref{eq:detpost}), we have
\begin{align}
\log\det{(\Gamma_m)}-\log\det{(\Gamma_{\rm post})}&=\log\det{(\Gamma_m)}-\log\det{\left(\frac{\det{(\Gamma_m)}}{(1+\sigma_e^{-2}w^\top\Gamma_m w)}\right)}\nonumber\\
&=\log\det{(\Gamma_m)}-\log\det{(\Gamma_m)}+(1+\sigma_e^{-2}w^\top\Gamma_m w)\nonumber\\
&=\log(1+\sigma_e^{-2}w^\top\Gamma_m w)\label{eq:logdet}
\end{align}

Now, we play the same trick, and let 
\[w=V\beta,\quad \beta\in\mathbb{R}^n,\]
with $V$ the eigenvectors of $\Gamma_m$ (see (\ref{eq:PCA})) and $\|w\|_2=\|\beta\_2=1.$ Plugging this into (\ref{eq:logdet}) we have 
\begin{align}
   \log\det{(\Gamma_m)}-\log\det{(\Gamma_{\rm post})}&=\log(1+\sigma_e^{-2}\beta^\top V^\top\Gamma_m V\beta)\nonumber\\
   &=\log(1+\sigma_e^{-2}\beta^\top\Lambda\beta)\nonumber\\
   &=\log(1+\sigma_e^{-2}\sum_{i=1}^n\lambda_i\beta_i^2),\nonumber
\end{align}
which (since $\beta_i^2\geq0$ and $\lambda_1\geq\lambda_2\geq\cdots\geq\lambda_n>0$) is again maximised when $\beta=e_1=[1,0,0,\dots,0]^T\in\mathbb{R}^n$. That is to say, setting $w=v_1$ reduces the the uncertainty the most (using D-opt).


\section{Learnability of linear prediction functionals}
Let $\CM \defeq \operatorname{Ran}(\prcov^{1/2})$
denote the Cameron--Martin space associated with the prior. 
This space is equipped with the
inner product
\[
    \ip{h_1}{h_2}_{\CM}
    \defeq
    \ip{\prcov^{-1/2}h_1}{\prcov^{-1/2}h_2},
    \qquad h_1,h_2\in\CM.
\]
The restriction of $\ell$ to $\CM$ has the Riesz representer
\begin{equation}\label{equ:CMrep}
    r_\ell \defeq \RR^{-1}_\CM (\ell\mid_\CM) = \prcov v_\ell.
\end{equation}
We call $r_\ell$ the Cameron--Martin representer of $\ell$.
In particular, for every $h\in\CM$,
$\ell(h)=\ip{v_\ell}{h}=\ip{r_\ell}{h}_{\CM}$.

\begin{proposition}\label{prp:unlearnable}
The posterior law of $\ell$ coincides with its prior law if and only if
its Cameron--Martin representer belongs to $\ker(\F)$.
\end{proposition}
\begin{proof}
Let $r_\ell$ be the Cameron--Martin representer of $\ell$ as in~\eqref{equ:CMrep}.
Note that if $r_\ell\in\ker(\F)$, then 
by~\eqref{equ:moments}, the prior and posterior laws of $\ell$ have the same
mean and variance. Since both laws are Gaussian, they coincide.
Conversely, assume the posterior law of $\ell$ coincides with its prior
law. Thus, $\lprsig = \lpostsig$ and therefore 
$\ip{\F r_\ell}{\Gammay^{-1}(\F r_\ell)} = 0$.
Since $\Gammay$ is strictly positive, 
$\F r_\ell=0$.
\end{proof}
This result shows that a prediction functional remains entirely uninformed by data
when its Cameron--Martin representer lies in $\ker(\F)$. 
\begin{remark}
The condition in Proposition~\ref{prp:unlearnable} has a simple statistical
interpretation. Under the joint prior law of $(m,y)$, we have
\begin{equation}\label{equ:derivation}
\begin{aligned}
\Cov\bigl(y,\ell(m)\bigr)
&=
\mathbb{E}\!\left[
\bigl(y-\mathbb{E}[y]\bigr)
\bigl(\ell(m)-\mathbb{E}[\ell(m)]\bigr)
\right] \\
&=
\mathbb{E}\!\left[
\bigl(\F(m-\dparpr)+\eta\bigr)
\ip{v_\ell}{m-\dparpr}
\right] \\
&=
\F\mathbb{E}\!\left[
(m-\dparpr)\ip{v_\ell}{m-\dparpr}
\right]
+
\mathbb{E}\!\left[
\eta\ip{v_\ell}{m-\dparpr}
\right] \\
&=
\F\prcov v_\ell
=
\F r_\ell.
\end{aligned}
\end{equation}
Hence, $r_\ell\in\ker(\F)$ if and only if $\ell(m)$ is uncorrelated with the
observations. Since $(\ell(m),y)$ is jointly Gaussian, this is equivalent to
independence.  Thus, the posterior law of $\ell$ coincides with its prior law
precisely when $\ell(m)$ is independent of the data.

To justify the derivation in~\eqref{equ:derivation}
we note that the second term in the third line vanishes by the independence of $m$
and $\eta$, and the final equality follows from the definition of the prior
covariance operator. Namely,
for an arbitrary $w\in\M$, 
\[
\ip{
\mathbb{E}\!\left[
(m-\dparpr)\ip{v_\ell}{m-\dparpr}
\right]
}{w}
=
\mathbb{E}\!\left[
\ip{v_\ell}{m-\dparpr}
\ip{m-\dparpr}{w}
\right] 
=
\ip{\prcov v_\ell}{w}.
\]
Thus,
$\mathbb{E}[(m-\dparpr)\ip{v_\ell}{m-\dparpr}]
=
\prcov v_\ell$.
We have also used the fact that 
for the Bochner integrable $\M$-valued random variable 
$X \defeq (m-\dparpr)\ip{v_\ell}{m-\dparpr}$,
$\ip{\mathbb{E}[X]}{w} = \mathbb{E}\!\left[\ip{X}{w}\right]$.
\end{remark}

As a measure of learnability of a linear functional, we consider 
the relative variance reduction
\[
    \LL(\ell)
    \defeq
    \frac{\lprsig-\lpostsig}{\lprsig}.
\]
We will refer to this ratio as the \emph{learnability} ratio.
Proposition~\ref{prp:unlearnable} identifies linear 
prediction functionals for which $\LL(\ell) = 0$. 
More generally, we can quantify the learnability of 
linear functionals by considering their learnability ratio. 
We first record the following basic result: 
\begin{lemma}
For each $\ell \in \M'$,
the learnability ratio satisfies the following:
\begin{enumerate}
\item $\LL(\ell) \in [0, 1]$; and 
\item 
letting $\Gammay$ be as in~\eqref{equ:Gammay},
we have 
\begin{equation}\label{equ:LL}
\LL(\ell) =  
\frac{\ip{\F\prcov v_\ell}{\Gammay^{-1} (\F\prcov v_\ell)}}{\ip{\prcov v_\ell}{v_\ell}}. 
\end{equation}
\end{enumerate}
\end{lemma}
\begin{proof}
The result follows from the definition of the prior and posterior laws of 
$\ell$.
\end{proof}

We next consider the following operator, which plays a key role in the analysis that follows.
\begin{equation}\label{equ:St}
    \St \defeq \prcov^{1/2}\F^*\Gammay^{-1}\F\prcov^{1/2}.
\end{equation} 
The dominant eigenvectors of $\St$ determine a set of canonical learnable directions.
We first note the following basic result.
\begin{lemma}\label{lem:S-spectrum}
The following hold.
\begin{enumerate}
\item $\St$ is a bounded self-adjoint operator;
\item  $0\preceq \St \prec \I$;
\item  if $\F \neq 0$, then $\St \neq 0$.
\end{enumerate}
\end{lemma}
\begin{proof}
The boundedness and self-adjoint property are straightforward to note.
To show (2), we first note 
that 
\[
\St
=
\bigl(\Gammay^{-1/2}\F\prcov^{1/2}\bigr)^*
\bigl(\Gammay^{-1/2}\F\prcov^{1/2}\bigr)
\succeq 0.
\]
To see $\St \prec \I$, we factorize $\St$ as follows: 
let
$\B\defeq \noise^{-1/2}\F\prcov^{1/2}$ and note that
$\St = \B^*(\I+\B\B^*)^{-1}\B$.
Subsequently, the identity
\[
    \I-\B^*(\I+\B\B^*)^{-1}\B
    =
    (\I+\B^*\B)^{-1}
\]
yields
$\I-\St=(\I+\B^*\B)^{-1}\succ0$.
Hence, $0\preceq\St\prec \I$. 
To see (3), note that since 
$\Gammay \succ 0$, $\St = 0$ implies $\F\prcov^{1/2} = 0$.  
This along with 
density of $\operatorname{Ran}(\prcov^{1/2})$ and boundedness of $\F$ imply
$\F=0$. 
\end{proof}
Note that as long as $\F \neq 0$, which is an underlying assumption, we have $\St\neq0$. 
In particular, we have that $\lambda_{\max}(\St) \in (0, 1)$. 


\begin{proposition}\label{prp:max-learnability}
Among all nonzero bounded linear functionals
$\ell\in\M'$, the learnability ratio $\LL(\ell)$ is maximized by
$\ell^* \defeq \RR_{\M}\bigl(\prcov^{-1/2}w\bigr)$,
where $w$ is an eigenvector corresponding to the largest eigenvalue of
$\St$ in~\eqref{equ:St}. Also,
\[
    \max_{\ell\in\M'\setminus\{0\}}\LL(\ell)
    =
    \lambda_{\max}(\St).
\]
\end{proposition}
\begin{proof}
Note that $w = \lambda_{\max}(\St)^{-1}\St w \in \operatorname{Ran}(\prcov^{1/2})$.
Thus, $\prcov^{-1/2}w$ is well-defined.

Next, let $\ell\in\M'\setminus\{0\}$ and consider
\[
    v_\ell\defeq \RR_{\M}^{-1}\ell \quad\text{and}\quad
    r_\ell
    \defeq
    \RR_{\CM}^{-1}\bigl(\ell|_{\CM}\bigr) = \prcov v_\ell.
\]
Hence, 
letting 
$z_\ell\defeq \prcov^{-1/2}r_\ell = \prcov^{1/2}v_\ell$, 
it follows from~\eqref{equ:LL} that
\[
    \LL(\ell)
    =
    \frac{\ip{\St z_\ell}{z_\ell}}{\ip{z_\ell}{z_\ell}}
    \leq
    \lambda_{\max}(\St).
\]
Finally, taking
$\ell^*\defeq\RR_{\M}\bigl(\prcov^{-1/2}w\bigr)$
yields $z_\ell=w$, and thus
$\LL(\ell^*)=\lambda_{\max}(\St)$.
\end{proof}

\begin{remark}
Let $\{w_j\}_{j\geq 1}$ denote orthonormal eigenvectors of $\St$
corresponding to its positive eigenvalues $\{\lambda_j\}_{j\geq 1}$.
Since
$w_j = \lambda_j^{-1}\St w_j \in \operatorname{Ran}(\prcov^{1/2})$,
the vectors
$v_j \defeq \prcov^{-1/2}w_j$
are well-defined. The associated functionals
\[
    \ell_j(\dpar) \defeq \ip{v_j}{\dpar}
\]
may be viewed as canonical prediction functionals for the inverse problem.
Their learnability ratio satisfy
$\LL(\ell_j)=\lambda_j$.
Hence, they are ordered according to the relative reduction in their prior
variances induced by the observations.
\end{remark}

In practice, one does not get to choose an arbitrary prediction 
functional aligned with eigenvectors of $\St$. However, one might 
have to choose between several physically meaningful prediction/goal 
functionals. The following result provides a quantitative estimate:
\begin{proposition}\label{prop:learnability-bound}
Let $\ell\in\M'\setminus\{0\}$, and let $r_\ell$ denote its
Cameron--Martin representer. Then the following hold.
\begin{enumerate}
\item
For every $\epsilon\geq0$,
$\LL(\ell)\leq\epsilon$
if and only if
$\|\Gammay^{-1/2}\F r_\ell\| \leq \epsilon^{1/2}\|r_\ell\|_{\CM}$.
\item
We have $\LL(\ell)\leq\epsilon$ whenever
$\|\noise^{-1/2}\F r_\ell\| \leq \epsilon^{1/2}\|r_\ell\|_{\CM}$.
\end{enumerate}
\end{proposition}
\begin{proof}
To see (1), note that by the definition of $r_\ell$,
%\[
%    r_\ell=\prcov v_\ell, \qquad \|r_\ell\|_{\CM} = \|\prcov^{1/2}v_\ell\|.
%\]
%Hence,
\begin{equation}\label{equ:LLCM}
    \LL(\ell)
    =
    \frac{\|\Gammay^{-1/2}\F r_\ell\|^2}{\|r_\ell\|_{\CM}^2},
\end{equation}
Further, we note that $\Gammay^{-1}\preceq\noise^{-1}$.
Therefore,
$\|\Gammay^{-1/2}\F r_\ell\| \leq \|\noise^{-1/2}\F r_\ell\|$.
Thus, (2) follows from (1).
\end{proof}
The first statement of the above result 
provides an exact characterization of the learnability ratio in
terms of the Cameron--Martin representer $r_\ell$. In particular, the quantity
in the right hand side of~\eqref{equ:LLCM}
may be used to assess the learnability
of a prescribed prediction functional. Moreover,  
the second statement of the proposition gives a sufficient condition involving the noise covariance alone.
Thus, if the whitened image $\noise^{-1/2}\F r_\ell$ is small relative
to the Cameron--Martin norm of $r_\ell$, then $\ell$ is poorly learnable.

\section{Conclusion}
Insert conclusions here.



\section{One Dimensional Example}
Suppose 
\[\obs=\langle g,m\rangle=\B u +e,\]

with $u$ satisfying
\begin{equation}
\begin{aligned}\label{eq: fwd}
-\frac{d}{dx}(\kappa(x)\frac{d}{dx}u)&=m\quad \text{in } (0,1)\\
u(0)&=u(1)=0,
\end{aligned}
\end{equation}
with $\kappa>0$ up to us. Assume we measure in an interval (ball) of radius $r>0$ so that 

\[\B u:=\int_{x_0-r}^{x_0+r} u(x) \;dx\]

Then we can write the measurement (in terms of $m$) as

\[\B u= \int_{x_0-r}^{x_0+r}  \int_{0}^{1} G(s,x;\kappa)m \;dx\;ds=  \int_{0}^{1} \int_{x_0-r}^{x_0+r} G(s,x;\kappa)m \;ds\;dx=\langle g_\kappa,m\rangle,\]
with $G(s,x;\kappa)$ the Green's function related to (\ref{eq: fwd}), and $g_\kappa=\int_{x_0-r}^{x_0+r} G(s,x;\kappa)\;ds$. Now, if $\mathcal{C}=\sum_{i=1}^\infty\lambda_i\phi_i\otimes\phi_i$,
we want to try to select $\kappa$ such that $g_\kappa=\phi_1$, that is pick $\kappa$ such that $\int_{x_0-r}^{x_0+r} G(s,x;\kappa)\;ds=\phi_1$. Clearly $\phi_1$ is dependent on the choice of $\mathcal{C}$. As a concrete example, suppose $\mathcal{C}=\mathcal{A}^{-2}$ with $\mathcal{A}:=\alpha I -\Delta$ with homogeneous Dirichlet boundary conditions. Then $\phi_1=\sin(\pi x)$.

\bibliographystyle{plain}
\bibliography{main}
\end{document}